% Chapter 56: Nuclei and Radioactivity (Supplement)
\chapter{Nuclei and Radioactivity (Supplement)}

Chapter~51's \(\alpha\)-particle gold-foil scattering put nearly all atomic mass in a nucleus a hundred-thousandth the atom's size. Five chapters treated it as an unchanging positive ball, binding electron clouds and determining chemistry and periodicity. Now enter that ball. It is not static hardware but a soup that can reshuffle spontaneously, changing identity and emitting energetic fragments: radioactivity, alchemists' dreamed-of elemental transmutation, driven by probability rather than sorcery.

\step{A Counterintuitive Opening}

Your body is a radiation source. Of roughly 140 grams of potassium, a tiny fraction is unstable potassium-40, undergoing about four thousand decays per second and emitting fast electrons or energetic photons into tissue. Carbon-14 in bones and fat adds roughly three thousand. Your body hosts seven or eight thousand miniature nuclear explosions every second, over two quadrillion yearly, without visible glow, abnormal warmth, or sensation.

% BEGIN TRANSLATOR SAFETY NOTE
\textbf{Translator safety note:} The uranium-salt manipulations below are historical examples, not experiments to repeat. Do not grind, heat, dissolve, or otherwise process radioactive materials; dust and solutions can create internal exposure and chemical hazards. See \href{https://www.epa.gov/radiation/radionuclide-basics-uranium}{EPA uranium health guidance}.
% END TRANSLATOR SAFETY NOTE

In 1896 Becquerel put uranium salts on wrapped photographic plates while studying fluorescence. Cloudy weather postponed his experiment, leaving salts in a drawer for days, yet the plates exposed precisely beneath them. Without light, electricity, or intervention, a rock continuously emitted something penetrating black paper. Curie later found emission independent of temperature, chemistry, or storage time: grind, dissolve, freeze, or heat salts red-hot and emission remains unchanged. An uncontrollable, unending mechanism hid in stone.

Another mystery hangs from your ceiling. Many smoke alarms contain a button-sized radioactive americium-241 source emitting \(\alpha\) particles continuously. Radioactivity is not confined to power stations; it lives in bodies, houses, and soil. Why and when do nuclei decay? Where does the radiation go? Why survive seven or eight thousand self-bombardments per second?

\step{Make a Guess First}

Write your intuitions before continuing.

First: positive protons repel inside a nucleus with roughly one-femtometer radius, \(10^{-15}\)~m. Later neutrons add no charge. What binds it: gravity, residual electromagnetic attraction, or a new interaction? Chapter~19's Coulomb law gives proton repulsion around 1.4 million electronvolts at one femtometer, with gravity unable to supply even a trillionth.

Second: how many decays per second from one gram of radium, hundreds, millions, or more? After a thousand years buried, how much weaker is its radiation?

Third: one hundred radioactive atoms have a one-day half-life, leaving fifty on average after a day, twenty-five after two, twelve or thirteen after three. When do all die? More fundamentally, can you predict whether one particular atom decays tomorrow? Does surviving one day make it tired and closer to death?

\step{Hands-on Experiment}

\begin{experiment}[Simulate Decay with One Hundred Coins]
\textbf{Materials}: fifty to two hundred coins, buttons, beans, or dice; paper and pencil.

\textbf{Step one}: count \(N_0\) coins and record the first table row. Scatter them on a flat surface and remove heads, the decayed ones. Count remaining \(N\) after round 1.

\textbf{Step two}: collect survivors, scatter again, remove heads. Repeat eight to ten rounds, recording survivors each time.

\textbf{Step three}: plot survivors against rounds and compare with exact halving from \(N_0\), through \(N_0/2\), \(N_0/4\), \(N_0/8\), and so on.

\textbf{Record and think}: points roughly follow the curve but fluctuate. Can you predict one coin's next removal? Why do a hundred coins behave more regularly than ten? After five surviving rounds, is a coin more likely to leave on round six? Save answers for comparison with nuclei in the opening explanation.
\end{experiment}

Actually do this, even with thirty coins: it unlocks our statistics. Each round is a fixed interval, each coin a nucleus, heads decay within that interval. Unpredictable individuals and orderly crowds reflect radioactivity's deepest property, not imperfect instruments: no nuclear countdown exists in principle.

\step{The Old Model Runs into Trouble}

Apply Chapter~19 electromagnetism and Chapter~52 atoms to nuclei and see where they fail.

First, nuclei should not exist. Two helium protons a femtometer apart have repulsion energy around \(1.4\times10^{6}\)~eV, requiring thermal temperatures above \(10^{10}\)~K, nearly a thousand solar-core \(1.5\times10^{7}\)~K. Gravity supplies only roughly a \(10^{-36}\) fraction of that energy, negligible. Like a handful of mutually repelling magnets, nuclei should explode under trillion-degree-equivalent repulsion, yet uranium sits in rocks for billions of years. The missing ingredient is an entire interaction, not a correction.

Second, mass and charge disagree. Helium has four hydrogen-nucleus masses but two charges; carbon twelve masses and six charges. Proton-only nuclei would keep proportionality. A historical patch added internal electrons to neutralize surplus protons without much mass. But Chapter~44's uncertainty confines electrons at femtometer scales only with tens-of-MeV momentum energies, too high for nuclei to retain such light particles; measured nuclear spins also disagree. Chadwick buried the patch in 1932, finding neutral, proton-mass neutrons in beryllium radiation. Two protons and two neutrons give helium's four masses and two charges.

Third, \(\beta\) decay seems to violate energy conservation. An emitted electron, \(\beta\) radiation, raises nuclear charge by one. Identical parent-to-daughter transformations should emit fixed electron energies, like Chapter~51's \(\alpha\) energies. Instead electrons span continuously from near zero to an endpoint, with recoil momenta also unbalanced. Either conservation fails or an unseen third party takes the remainder. Pauli chose the latter in 1930: a neutral, almost noninteracting particle secretly carries missing energy and momentum. Named the neutrino, it was directly detected twenty-six years later. Conservation survived; the visible-particles-only account failed.

\step{Inventing New Concepts}

Three failures point to one blueprint. Invent the needed pieces.

First, \textbf{nucleons}, protons and neutrons. Their nuclear roles are nearly equal, masses differing by under 1.4 parts per thousand, neither afraid of mutual electric repulsion, neutrons having none. Proton count \(Z\) determines element, neutral-electron count, and chemistry; neutron count \(N\) completes identity, with mass number \(A=Z+N\). Same protons, different neutrons make \textbf{isotopes}, chemically indistinguishable but nuclear opposites. Hydrogen, deuterium, and tritium are brothers, two stable and one radioactive.

Second, \textbf{the strong interaction}. A new attraction overwhelms Coulomb repulsion near one femtometer but almost vanishes beyond two or three, absent from ordinary chemistry, syringes, and tables. It is roughly charge-independent across proton/proton, neutron/neutron, and proton/neutron pairs. Saturation makes each nucleon hold only a few immediate neighbors, unlike long-range charge interactions. These three traits shape nuclei.

Third, \textbf{binding energy and mass defect}. Separating nucleons costs binding energy \(B\). Chapter~40's mass--energy relation says supplied energy adds mass, so a bound nucleus's rest mass \(M\) is less than its free constituents:
\[
B=\bigl(Zm_{\mathrm p}+Nm_{\mathrm n}-M\bigr)c^{2}.
\]
The assembled soup weighs less than ingredients separately; missing mass records binding energy. This is not mystical mass conversion but direct evidence that energy has inertia and appears on a scale. Mass here always means intrinsic invariant rest mass.

Fourth, \textbf{three decays}. Unstable nuclei seek steadier identities. \textbf{\(\alpha\) decay} ejects helium-4, two protons and two neutrons, reducing \(A\) by four and \(Z\) by two, changing element. \textbf{\(\beta\) decay} converts neutron to proton or vice versa while emitting an electron or positron and a neutrino; \(A\) stays, \(Z\) changes by one. This identity change belongs to the fourth fundamental interaction, weak, alongside gravity, electromagnetism, and strong. \textbf{\(\gamma\) decay} drops a nucleus's energy, emitting an energetic photon without changing composition. Penetration differs: paper or surface skin stops \(\alpha\), millimeters of aluminum stop \(\beta\), thick lead or concrete stop \(\gamma\), distinctions vital for safety later.

Fifth, our heart: \textbf{decay constant and half-life}. Each unstable nucleus has a fixed unit-time decay probability \(\lambda\), independent of age, survivor count, temperature, or pressure. This is an \emph{individual} probability, not a collective appointed death. Large groups therefore lose half on average every fixed \(T_{1/2}\), the half-life, a large-number statistical pattern. Individual nuclei have no age, record, or approaching deadline, as your coins anticipated.

\step{The Model in One Sentence}

\begin{oneline}
Short-range strong interactions bind protons and neutrons, most tightly per nucleon near iron. Less tightly bound nuclei randomly change identity at fixed probability, emitting \(\alpha\), \(\beta\), or \(\gamma\); populations decay exponentially by half-lives while individual times remain unpredictable.
\end{oneline}

Like a crowded waiting room where each passenger repeatedly rolls a private die and leaves on one result, with equal constant chances. Larger crowds make total departures predictable, never one named person's round.

Three limits: passengers tire and grow more likely to leave, but billion-year-old and newborn uranium nuclei have equal next-second probability, memorylessness. Human departures have triggers; quantum nuclear decay offers no deeper observable trigger, sharing Chapters~44 and~45's tunneling and probability. Finally, isotope and \(\lambda\) have no adjustable universal knob: strong, weak, and Coulomb competition fixes them, explaining Becquerel's unshutoffable rock.

\step{The Mathematical Translator}

Calculate helium-4's two protons and two neutrons using \(m_{\mathrm p}=1.007276\)~u, \(m_{\mathrm n}=1.008665\)~u, and measured nuclear \(M=4.002603\)~u, with \(1~\mathrm{u}\,c^{2}=931.5\)~MeV. Defect:
\[
\Delta m=2\times1.007276+2\times1.008665-4.002603=0.030377~\mathrm{u},
\]
\[
B=\Delta m\,c^{2}=0.030377\times931.5\approx28.3~\mathrm{MeV},
\qquad \frac{B}{A}\approx\frac{28.3}{4}\approx7.1~\mathrm{MeV}.
\]
Four questions: total and per-nucleon disassembly cost; mass deficit times \(c^2\) converts the account; exact for rest masses as experimental fact rather than a model; no internal structure or decay prediction follows. Check \(Z\) or \(N\) zero: \(B=0\), one proton has no binding. Negative deficit would make a nucleus heavier than its ingredients, so nature would not assemble it; nature makes only \(B>0\) nuclei.

Plot \(B/A\) across elements: climb from hydrogen, a helium-4 spike, peak near iron \(A\approx56\) around \(8.8\)~MeV, then descend slowly toward uranium's \(7.6\)~MeV. This nuclear terrain makes movements toward iron, light fusion or heavy fission, climb binding per nucleon while descending total energy and releasing it. Stars and nuclear stations use opposite slopes of one map.

\begin{figure}[htbp]
\centering
\begin{tikzpicture}[>=Stealth,scale=1.0]
  \draw[->] (0,0) -- (10.2,0) node[below left=1pt]{Nucleon count \(A\)};
  \draw[->] (0,0) -- (0,5.6) node[left=1pt]{\(B/A\) (MeV)};
  \foreach \y/\lab in {1/2,2/4,3/6,4/8}{
    \draw (-0.08,\y*0.6+0.6) -- (0.08,\y*0.6+0.6) node[left=4pt]{\scriptsize \lab};
  }
  % Binding per nucleon: hydrogen, rapid rise, Fe peak, slow decline
  \draw[thick,blue!60!black] plot[smooth] coordinates {
    (0.15,0.6) (0.5,3.9) (1.0,4.4) (2.0,4.9) (3.5,5.25) (4.5,5.3)
    (5.5,5.2) (7.0,4.95) (8.5,4.75) (9.7,4.6)};
  \fill (4.5,5.3) circle (1.5pt) node[above=2pt]{\scriptsize Iron \(A\approx56\)};
  \fill (0.5,3.9) circle (1.5pt) node[above right=0pt]{\scriptsize Helium-4};
  \fill (9.7,4.6) circle (1.5pt) node[above=1pt]{\scriptsize Uranium};
  \draw[->,red!60!black,thick] (1.2,5.5) -- (3.4,5.5) node[midway,above]{\scriptsize Fusion: light nuclei climb, release energy};
  \draw[->,red!60!black,thick] (9.0,5.9) -- (6.2,5.9) node[midway,above]{\scriptsize Fission: heavy nuclei climb, release energy};
\end{tikzpicture}
\caption{Binding energy per nucleon peaks near tightly bound iron. Both light-nucleus fusion and heavy-nucleus fission move toward iron and release energy.}
\end{figure}

Now decay statistics. Starting with \(N_0\) identical radioactive nuclei, after \(t\) the average is
\[
N(t)=N_{0}\times2^{-t/T_{1/2}} .
\]
Check \(t=0\), \(N_0\); \(t=T_{1/2}\), \(N_0/2\); \(t=2T_{1/2}\), \(N_0/4\). As \(t\) tends to infinity, \(N\) approaches but never reaches zero. Thus average extinction never finishes, only becomes invisible; one hundred atoms have no predetermined final day.

\begin{mathbridge}
Why exponential? Constant per-nucleus rate \(\lambda\) means a short \(dt\) gives mean \(\lambda N\,dt\) decays:
\[
\frac{dN}{dt}=-\lambda N .
\]
The key \(N\) makes fewer survivors decay more slowly. Solution:
\[
N(t)=N_{0}e^{-\lambda t}.
\]
Set \(N=N_0/2\) to get \(T_{1/2}=\ln2/\lambda\approx0.693/\lambda\). Factors \(2^{-t/T_{1/2}}\) and \(e^{-\lambda t}\) differ only in base. Surviving \(s\) then another \(t\) has conditional probability \(e^{-\lambda(s+t)}/e^{-\lambda s}=e^{-\lambda t}\), independent of \(s\). Waiting changes nothing; probability starts afresh.
\end{mathbridge}

Activity \(A=\lambda N\) counts decays per second in becquerels, one Bq per second. A seventy-kilogram adult contains about \(140\)~g potassium; natural potassium-40 fraction \(0.0117\%\) gives \(0.0164\)~g, or \(N=0.0164/40\times6.0\times10^{23}\approx2.5\times10^{20}\) atoms. Half-life \(T_{1/2}=1.25\times10^{9}\) years is about \(3.9\times10^{16}\)~s:
\[
A=\lambda N=\frac{\ln2}{T_{1/2}}N\approx\frac{0.693\times2.5\times10^{20}}{3.9\times10^{16}}\approx4.4\times10^{3}~\mathrm{Bq}.
\]
Four thousand four hundred per second matches the opening. One gram radium-226 has \(N=6.0\times10^{23}/226\approx2.7\times10^{21}\) atoms, half-life \(1600\) years, and activity \(3.7\times10^{10}\)~Bq, historically defining the curie. Guess two thus gives thirty-seven billion per second, four orders beyond millions, and about two-thirds strength after a thousand years, roughly 0.6 half-lives. Nuclear numbers outstrip intuition.

Activity is not injury. Safety distinguishes Bq decays, absorbed dose Gy in joules per tissue kilogram, and equivalent dose Sv weighted by radiation type. Equal deposited energy gives \(\alpha\) roughly twenty times \(\gamma\)'s biological damage, hence weighting. Natural cosmic, soil, radon, and internal-potassium background averages \(2.4\)~mSv yearly; chest X-ray about \(0.1\)~mSv, chest CT \(7\)~mSv. Below annual \(100\)~mSv, epidemiology detects no distinguishable added risk. Confusing these accounts fuels most radiation panic.

% BEGIN TRANSLATOR SAFETY NOTE
\textbf{Translator safety note:} Do not treat 100 mSv per year as a safe allowance or a threshold below which risk is absent. Low-dose ionizing radiation can increase long-term cancer risk; an individual's risk depends on exposure and other factors. Discuss medical imaging's benefits and risks with the clinician, rather than using this paragraph to seek or avoid a scan. See \href{https://www.who.int/news-room/fact-sheets/detail/ionizing-radiation-and-health-effects}{WHO radiation-health guidance}.
% END TRANSLATOR SAFETY NOTE

Energy next: uranium-235 fission releases about \(200\)~MeV, \(3.2\times10^{-11}\)~J. A gram's \(2.6\times10^{21}\) nuclei yield \(8.2\times10^{10}\)~J, roughly \(2.8\) tonnes of standard coal or over ten household electricity years. One gram for nearly three tonnes is binding-curve mass--energy exchange. The Sun radiates \(3.8\times10^{26}\)~J each second; \(E=mc^2\) corresponds to about \(420\) times ten thousand tonnes of rest mass, continuously for 4.6 billion years through hydrogen-to-helium ascent.

\step{The Model's Limits}

The femtometer strong-attraction picture is effective, not a simple spring. Its distance dependence rapidly vanishes beyond two femtometers and reverses to repulsion very close, preventing point collapse. Deeper, composite protons and neutrons contain quarks; nuclear force is residual quark/gluon interaction, like intermolecular forces are residual electromagnetism. Chapter~59 expands that map. Do not extrapolate our effective force beyond its scale.

Binding curves are terrain, not universal formulas. Shell structure ripples atop the trend. Nuclear proton and neutron shells are especially stable at these magic counts: \(2,8,20,28,50,82,126\). Specific decay channels require shell details, Coulomb effects, and weak interactions, not the smooth curve alone.

Remember half-life's statistical boundary. Billions give precise half-populations; ten give a fluctuating average, as coins show. Individuals have probabilities, not half-lives. Exponentials also fail at extremes: few survivors are fluctuation-dominated, and intervals below \(10^{-23}\) seconds may precede exponential decay. These irrelevant-to-daily-lab corrections nevertheless mark a model, not dogma.

The \(\beta\) spectrum teaches methodology. Apparent failed conservation actually missed an almost invisible particle. An unbalanced account prompts searching omitted terms, not declaring laws dead. Yet neutrinos were established by real detection twenty-six years later, not accounting cleverness alone. Hypotheses must become measurements.

Two radiation myths need correction. Irradiated objects do not automatically become sources: activation requires neutron bombardment changing nuclei. X-ray or \(\gamma\)-irradiated patients and food retain unchanged nuclei and do not themselves radiate afterward; irradiation's radioactivity leaves with its rays. Nor does every ever-smaller dose have measurable harm. Below about \(100\)~mSv yearly, added risk cannot be distinguished from present statistical noise; linear-no-threshold extrapolation is cautious management, not established fact. Treating it as fact breeds banana and CT panic. Natural radiation's largest single source is ground-emitted indoor radon, not nuclear plants or medicine; ventilation is an ordinary person's most practical action.

% BEGIN TRANSLATOR SAFETY NOTE
\textbf{Translator safety note:} The repeated low-dose reassurance is not a safety threshold: \href{https://www.iarc.who.int/pressrelease/low-doses-of-ionizing-radiation-increase-risk-of-death-from-solid-cancers/}{IARC reports increased cancer risk with protracted low-dose exposure}. For home radon, do not rely on ventilation alone; test and obtain appropriate mitigation if elevated. See \href{https://www.epa.gov/radtown/radon-homes-schools-and-buildings}{EPA radon testing and reduction guidance}.
% END TRANSLATOR SAFETY NOTE

\begin{challenge}
Why half-lives from uranium-238's \(45\) times a hundred million years to polonium-212's \(0.3\) microseconds, a \(10^{24}\) ratio, when \(\alpha\) energies differ only from \(4.3\) to \(8.8\)~MeV? Gamow's 1928 answer applies Chapter~44 tunneling: \(\alpha\) particles trapped by nuclear attraction face a Coulomb barrier, classically impassable but quantum-mechanically penetrable with \(P\approx e^{-2G}\). The quantity \(G\) is highly barrier-sensitive. Higher \(\alpha\) energy lowers and narrows the effective barrier, exponentially increasing escape. Doubling energy can change probability by over twenty orders, joining billions of years and microseconds. Thus \(\alpha\) decay was an early everyday manifestation of tunneling.
\end{challenge}

\step{Explaining the Opening Phenomenon}

Retrieve all three opening stories.

Your seven or eight thousand decays per second come from natural potassium-40 and carbon-14. Why unavoidable? Kidneys and membranes cannot distinguish potassium-39 and potassium-40 with identical electrons; dietary potassium always includes its radioactive fraction. This is Earth's geochemical original equipment, not pollution: potassium-40 existed in Earth's four-billion-year-old ingredients and outlasts calmly with a 1.25-billion-year half-life. Life has always occupied this background. Tiny decay energies spread among tens of trillions of cells, whose evolved repair machinery readily handles annual internal doses around one millisievert. No radiation-free environment ever offered an alternative.

% BEGIN TRANSLATOR SAFETY NOTE
\textbf{Translator safety note:} Natural origin and cellular repair do not guarantee harmlessness or justify additional exposure. Radiation risk depends on dose and exposure conditions; see \href{https://www.who.int/news-room/fact-sheets/detail/ionizing-radiation-and-health-effects}{WHO radiation-health guidance}.
% END TRANSLATOR SAFETY NOTE

Becquerel's salts cannot switch off because \(\lambda\) belongs to nuclear strong, weak, and Coulomb balance, while heat, pressure, and chemistry affect electron clouds at energies hundreds of thousands of times smaller. Fire or acid perturbs electrons without the nuclear soup noticing. Uranium-238 and descendants' \(\beta\) and \(\gamma\) penetrated black paper in 1896. Its 4.5-billion-year half-life leaves almost unchanged emission today.

% BEGIN TRANSLATOR SAFETY NOTE
\textbf{Translator safety note:} This comparison is not an experiment to repeat: do not heat, dissolve, grind, or otherwise process uranium or other radioactive samples. Internal contamination and chemical toxicity are additional hazards; see \href{https://www.epa.gov/radiation/radionuclide-basics-uranium}{EPA uranium guidance}.
% END TRANSLATOR SAFETY NOTE

Smoke alarms domesticate radioactivity. Americium's \(\alpha\) ionizes chamber air to maintain weak steady current; smoke captures ions, reducing current and triggering sound. Bedroom use is acceptable because \(\alpha\) cannot cross centimeters of air and a plastic housing; unless dismantled and swallowed, dose is negligible. The same paper-stopped external radiation deposits all its energy next to cells if inhaled or ingested, explaining radon's asymmetry. Source location relative to the protective wall matters.

% BEGIN TRANSLATOR SAFETY NOTE
\textbf{Translator safety note:} Keep smoke alarms installed and use them as directed. Never dismantle the ionization chamber, remove its source, or handle radioactive components; follow manufacturer and local disposal instructions. Approved intact alarms save lives and have negligible radiation risk. See \href{https://www.nrc.gov/reading-rm/doc-collections/fact-sheets/smoke-detectors}{NRC smoke-detector guidance}.
% END TRANSLATOR SAFETY NOTE

Fission and fusion climb opposite slopes. Uranium-235 absorbs a neutron and splits into closer-to-iron fragments, turning increased binding into kinetic energy and releasing two or three neutrons for a chain reaction. Moderators, control rods, and containment manage the neutron flow. Fusion joins hydrogen isotopes into helium, climbing toward 7.1 MeV per nucleon of binding and releasing more. Positive nuclei must approach within a femtometer despite Coulomb repulsion; even the Sun's fifteen million degrees cannot classically surmount it. Chapter~44 tunneling lets protons penetrate instead, pacing billions of years of burning rather than an explosion. Terrestrial tokamaks and laser fusion struggle to suspend hundred-million-degree soup. Sun, reactors, alarms, and internal potassium are four footnotes to one nuclear terrain.

\step{Thought Experiments and Challenges}

\begin{thought}[A Strong Force with Ten Times the Range]
Hold everything else fixed and extend effective nuclear range from one to ten femtometers. Coulomb repulsion accumulates long-range, while strong attraction originally reaches only neighbors. Tenfold range gives each nucleon hundreds rather than a few partners, stabilizing heavy nuclei and shifting the iron peak much heavier. Uranium might cease to be the heaviest natural element, with stable superheavy crustal elements absent today. Conversely, halve strength or shorten range to a tenth of a femtometer and even helium-4 might fail; stars cannot ignite and hydrogen nearly monopolizes the universe. Carbon-life parameters have a startlingly narrow window: stronger or weaker rewrites periodicity. No available knob tests this counterfactual, but it reveals our element list as a delicate interaction balance, not a default entitlement.
\end{thought}

\begin{puzzles}
\item A million radioactive atoms have one-hour half-life. When is the mean below one, and what does that mean? Does half-life still matter? (Hint: find when \(2^n\) exceeds a million; \(2^{20}\approx10^{6}\). Near unit counts, fluctuations dominate and the large-number premise fails, leaving individual probability.)
\item Why use eighty-eight-year plutonium-238 for spacecraft heat rather than seven-hundred-million-year uranium-235 or 138-day polonium-210? Begin with power proportional to \(\lambda N\). (Hint: power \(=\lambda N\times\) energy per decay. Uranium's \(\lambda\) is too small for useful equal-mass power; polonium's large \(\lambda\) gives initial power but loses most within months, not ten-year missions. Eighty-eight years balances power and decades of endurance.)
\item Living organisms exchange carbon with atmospheric proportions; death stops exchange and carbon-14 decays with a 5730-year half-life. Why date three-thousand-year wood but not sixty-five-million-year dinosaur fossils? (Hint: half a half-life leaves about 70 percent, ample signal. Over ten thousand half-lives leaves \(2^{-10000}\), not even one atom; use longer-lived uranium--lead clocks instead.)
\item Why are \(\alpha\) decay's \(\alpha\) energies usually fixed while \(\beta\) electrons form a continuum, historically challenging conservation? (Hint: \(\alpha\) gives daughter plus \(\alpha\), two bodies whose energies conservation fixes. \(\beta\) gives daughter, electron, and unseen neutrino, sharing a fixed total among three. Omit the third and accounts never balance.)
\end{puzzles}

Three threads remain. Chapter~59 maps \(\beta\)'s weak interaction alongside the other forces and asks why their different appearances might share language. Chapter~58 derives repeatedly used energy and momentum conservation from time and space symmetry, not accidental experimental summaries. And Chapter~45's quantum probability supplies decay chances without schedules, a theoretical structure rather than measurement failure. Every nuclear-soup reshuffle uses the same deck of rules.
